% 第2章 交易的信息论基础与 AI 的作用位置
\section{交易的信息论基础与 AI 的作用位置}\label{sec:foundations}

在谈任何模型之前，得先把交易这件事本身说清楚。这一章不介绍算法，只回答一个问题：在一个由人和机器共同参与的市场里，凭信息赚钱这件事，数学上被允许到什么程度？只有把这个边界画出来，后面每一族方法的能力与失效条件才有参照。

\subsection{有效市场之争留下的，不是一个“是不是有效”的答案}

Fama（1970）\cite{fama1970} 把有效市场分成弱式、半强式、强式三档，本质是把“价格是否已反映某类信息”变成一个可检验命题。半个世纪过去，这场争论真正留下的结论不是“市场有效”，而是一个更微妙的东西：\textbf{价格对新信息的反应速度快到让大多数慢半拍的人赚不到钱，但快的人能赚到——而快是有成本的。}

Grossman 与 Stiglitz（1980）\cite{grossmanstiglitz1980} 的悖论把这个微妙之处变成了严格的逻辑：如果价格已经充分反映所有信息，那么没有人有动力去花费成本获取信息；可一旦没人去获取，价格又不可能充分反映信息。于是市场只能停在一个不动点上——价格足够有效率，使得获取信息的边际成本恰好等于边际收益，而\textbf{这个均衡里必然留着一部分未被消除的定价偏差}，它就是所有主动交易的利润来源。Shleifer 与 Vishny（1997）\cite{shleifer1997} 进一步指出，即便偏差被人看见，套利者也会因为资金约束与噪声交易者风险而无法把它立刻消除，这正是偏差能长期存在的原因。这个结论对本文至关重要：它意味着“AI 能不能赚钱”这个问题问错了，正确的问题是“AI 相对其他参与者，在信息获取的边际上是不是更便宜、更快、或者更强”。

\subsection{信息如何进入价格：微观结构的视角}

Kyle（1985）\cite{kyle1985} 的模型给出了信息进入价格的最简机制。市场里有一类知情交易者掌握私有信息，一类噪声交易者随机下单，做市商只能看到净订单流并从中间接推断信息。均衡下，价格以某个速率（即 Kyle 的 $\lambda$，价格冲击系数）向真实价值收敛：
\begin{equation}
  \lambda = \frac{1}{2}\sqrt{\frac{\Sigma_0}{\sigma_u^2}},
\end{equation}
其中 $\Sigma_0$ 是私有信息的先验方差，$\sigma_u^2$ 是噪声交易的方差。这个式子说的是：\textbf{知情交易者越多、信息越有价值，价格对订单流就越敏感，单位信息能换取的利润就越薄。}这是一种自动的“信息租金耗散”机制——它和 Grossman--Stiglitz 讲的是同一件事，只是从订单流的层面重新讲了一遍。

对 AI 交易而言，这条机制的含义是：任何被公开、被广泛使用的信号，其毛利都会被 $\lambda$ 自动压薄。这解释了为什么大量“因子被发表后收益衰减”的经验现象（McLean \& Pontiff 2016\cite{mcleanpontiff2016}）不是意外，而是结构性必然。也解释了为什么纯“回测漂亮”的策略在实盘上失效——回测里的 $\lambda$ 往往被设成一个常数甚至零。

\subsection{低信噪比：金融预测与其它机器学习任务的根本差别}

把机器学习搬到金融上，最先撞上的不是算法问题，而是信噪比问题。图像分类里，像素与标签之间有相当强的确定性关系；而在月度股票收益预测里，最强的模型能解释的方差也只有千分之几。

Gu、Kelly 与 Xiu（2020）\cite{gukellyxiu2020} 用一整套机器学习方法（含树模型与神经网络）预测美股横截面收益，报告的样本外 $R^2$ 量级只有\textbf{约 0.3\%--0.4\%}（月度）。这个数字单看小得可笑，但换算成年化多空组合的收益与夏普，已经相当可观。它揭示的道理是：\textbf{金融里“解释力低”和“能赚钱”并不矛盾，因为杠杆和横截面聚合可以放大微弱的信号；但反过来，“样本内 $R^2$ 高”也几乎不能说明任何问题，因为任何一个过参数化模型都能把噪声拟合得很漂亮。}

低信噪比带来三个连锁后果。其一，过拟合的风险被放到最大，任何在样本内调参的自由度都是一种“数据窥探”（White 2000\cite{white2000}）。其二，模型比较的统计功效很低，两个模型的 $R^2$ 差异常常落在噪声里。其三，信噪比低意味着\textbf{模型的边际信息增量必须以“经济上可交易”来检验}，而不是以统计指标来检验——这正是评测污染审计（第 \ref{sec:audit} 章）要处理的问题。

\subsection{从预测到收益：摩擦如何把缝隙削窄}

即便一个信号真的抓到了信息，它离“赚到钱”还隔着两段：组合构建与交易执行。

组合构建这一端，Markowitz（1952）\cite{markowitz1952} 的均值--方差框架给出了最优权重：
\begin{equation}
  w^\star = \arg\max_{w}\ \Big( w^\top \mu - \frac{\gamma}{2} w^\top \Sigma w \Big),
\end{equation}
其中 $\mu$ 是预期收益，$\Sigma$ 是协方差，$\gamma$ 是风险厌恶。问题在于，$\mu$ 的估计误差会被 $\Sigma^{-1}$ 放大——低信噪比地区，估计误差主导结论，这就是“误差最大化”（error maximization）的经典困境。机器学习对组合构建的贡献，很大程度上是在更稳健地估计 $\mu$（和 $\Sigma$）。

执行这一端，Almgren 与 Chriss（2000）\cite{almgrenchriss2000} 把“想卖一大笔”分解成两个互相拉扯的成本：等待太慢会让价格漂移（时机风险），抛得太快会砸盘（冲击成本）。最优执行路径是二者的折中，总成本是一支关于成交速度的 U 形曲线。这条路后来被强化学习重新打开（第 \ref{sec:rl} 章）。

把两端合起来，就能给“可交易缝隙”（本文工作术语，其与既有概念的对应见第 \ref{sec:intro} 章 §1.1 的对齐表）一个可操作的定义。

\begin{definition}[可交易缝隙]
设策略 $\pi$ 在风险预算下产生预期超额收益 $\mathbb{E}[r_\pi]$、交易成本与冲击成本 $C(\pi)$、以及与所承担风险相称的补偿 $\lambda_\pi$。称该策略存在\emph{可交易缝隙}，当且仅当扣除摩擦后的净风险调整收益
\begin{equation}
  G(\pi) = \mathbb{E}[r_\pi] - C(\pi) - \lambda_\pi
\end{equation}
在\textbf{多重检验校正后}仍然显著大于零，且在样本外与跨市场不反号。
\end{definition}

这个定义里每一个词都有分量。$C(\pi)$ 不是可以事后忽略的尾项，它常常是决定生死的那一项（Martin \& Sch\"oneborn 2011 的说法很直白：均值回归能赚，但要付成本\cite{martin2011}）。“多重检验校正后”是因为在一个你试验了一万次的策略空间里，总能挑出一条回测最好的——“显著”必须相对于你试过多少次而言（Harvey、Liu \& Zhu 2016\cite{harveylz2016}）。“样本外不反号”是最后的诚实检验。

\subsection{缝隙存在的一个必要条件}\label{sec:gapcondition}

把可交易缝隙再推一步，可以得到一条能解释很多现象的必要条件。设某信号在 $t$ 时刻的残余预测力（可转化为毛收益的部分）按 $t$ 指数衰减，幅度为
\begin{equation}
  \delta(t) = \delta_0\, e^{-\lambda t},
\end{equation}
其中 $\lambda$ 是信息被市场吸收的速率（本章微观结构一节的 Kyle $\lambda$ 决定了它）；记单次调仓的比例成本为 $c$。若按固定周期 $\Delta$ 调仓，则单位时间的净收益是毛收益减成本再除以周期：
\begin{equation}
  g(\Delta) = \frac{\delta_0\, e^{-\lambda\Delta} - c}{\Delta}.
\end{equation}
令 $x=\lambda\Delta$，一阶条件整理为
\begin{equation}
  (1+x)\,e^{-x} = \frac{c}{\delta_0} =: \rho .
\end{equation}
而 $(1+x)e^{-x}$ 在 $x\ge 0$ 上从 1 单调递减，因此：

\begin{proposition}[缝隙存在的必要条件]
在上述模型下，存在使 $g(\Delta)>0$ 的调仓周期，当且仅当 $\delta_0 > c$；最优调仓周期为 $\Delta^\star = x^\star/\lambda$，其中 $x^\star$ 是方程 $(1+x)e^{-x}=\rho$ 的正根。
\end{proposition}

这个命题只有两行推导，却把三件事绑在了一起：\textbf{信号的初始幅度 $\delta_0$ 必须超过交易成本 $c$，否则再好的模型、再聪明的调仓节奏都无力回天；最优调仓频率随信息衰减速度 $\lambda$ 上升而上升（衰减越快越要抢在信息消失前交易）；随成本占比 $\rho$ 上升而下降（越贵越要少动）。}这解释了为什么高频策略对成本极度敏感、也解释了为什么低频信号在成本上升时会被迫降频甚至停摆。

不过，只有把判据里的两个量估出来，它才不是一个同义反复——否则“$\delta_0>c$”就等于“能赚钱的策略就能赚钱”。两者都可以用可核验的方式估：

\begin{itemize}
  \item \textbf{$c$ 的估计}：用交易记录直接算——每笔成交均价与决策时点参考价的差额，加上佣金与税费，再除以成交金额，即得该笔的比例成本；多笔取中位数，避免极端单笔带偏。
  \item \textbf{$\delta_0$ 的估计}：用“事件窗平均实现收益”作代理——在信号触发后的 $k$ 期内，按信号方向构建的多空组合的平均超额收益，取事件样本均值 $\hat\delta_0$。它衡量的是“信号在信息尚未衰减时能兑现多少”。
\end{itemize}

\textbf{估计误差如何影响判据}。$\hat\delta_0$ 的标准误大致随事件数 $m$ 与单事件收益的分散度 $\sigma_r$ 变化，量级约为 $\sigma_r/\sqrt{m}$；$k$ 越大，$\hat\delta_0$ 越被信息衰减压低，因此 $k$ 应取信息半衰期的量级。由于两个量都有误差，判据应当用\textbf{置信下界}而非点估计来判断：
\begin{equation}
  \hat\delta_0 - z_{\alpha}\,\widehat{\mathrm{SE}}(\hat\delta_0) \;>\; \hat c ,
\end{equation}
其中 $z_\alpha$ 取单侧分位（如 5\% 对应约 1.645）。这一改动看似保守，却是必要的：在事件数少、分散度大的情形下，用点估计判断会\textbf{系统性地高估}缝隙的存在性，而这正是大量小样本研究所犯的错误。

它还给出一个实用的诊断：当你发现某策略“样本内很好、实盘亏钱”时，先按上面的办法估 $\hat\delta_0$ 与 $\hat c$；若置信下界已经落在 $\hat c$ 之下，问题不在模型，而在缝隙根本不存在。

\subsection{AI 在这条链条上的三个作用位置}

用上面这套语言，AI 的应用可以归成三个位置，它们的能力上限各不相同：

\begin{enumerate}
  \item \textbf{表示学习（改变信息的天花板）}。把非结构化数据（新闻、财报、订单簿、卫星图、电话会录音）转成可用的数值特征。这一环直接决定“信息存量”能用出多少——文本与订单簿这类数据的信噪比未必高，但存量极大，且长期被低估。深度学习、图神经网络、大模型主要在这里发力。
  \item \textbf{偏差估计（把信号做成预测）}。在给定特征下估计未来收益的条件分布。这里的天花板由市场效率与 $\lambda$ 决定，AI 能做的通常是“更充分、更非线性地利用已有信息”，而非创造新信息。监督学习、序列模型主要在这里。
  \item \textbf{决策收割（把预测做成仓位）}。决定买多少、什么时候买、怎么执行、如何风控。这一环的天花板由摩擦决定，AI 的贡献是把摩擦降得更低、把动态决策做得更优。强化学习、最优执行、组合优化主要在这里。
\end{enumerate}

这三个位置不是并列的三选一，而是一条流水线：表示学习决定了信息能榨出多少，偏差估计决定了信号有多准，决策收割决定了信号能变现多少。一个方法在某篇论文里表现好，往往是因为它在其中某一环做得比对手好，而不是“AI 更强”。把这一点想清楚，后面读任何一篇“AI 交易”论文时，就不会被它的指标牵着走。

\subsection{一个同构性的看法}

如果跳出金融，会发现这三个作用位置和推荐系统、广告竞价、供应链优化里的“感知--预测--决策”分层是同一套结构。差别只在约束：金融这边，\textbf{信噪比更低、分布更非平稳、环境是对抗性的（你的对手会针对你的模型调整）、并且可用于学习的“样本”受限于有限的历史}。这四条约束里，前三条在其他领域可以被工程手段缓解，唯独第四条无法绕过——历史只有一份，而且它是由一个会因你的参与而改变的过程生成的。这也是为什么金融机器学习的“样本量”问题，本质上不是数据不够大，而是\textbf{有效的、独立的样本不够多}。第五章起的每一族方法，都会回到这条约束上。

最后落一句本章要记住的话：\textbf{AI 在金融里不是一台印钞机，它是一台更高效的“信息--预测--决策”转化装置；装置效率再高，也受制于喂进去的信息还剩多少、以及转化出来的东西够不够付摩擦的账。}
